\chapter{State of the Art}
\label{cap:state-of-the-art}

This chapter situates the platform within the existing landscape. It first
surveys recruitment software and the gaps it leaves for a small,
GDPR-constrained team; then reviews the role of artificial intelligence in
hiring --- CV parsing, candidate--job matching and conversational assessment;
introduces the language-proficiency framework adopted for assessment; and
finally summarises the regulatory context that shaped the data model. The
review motivates the design choices defended in Chapter~\ref{cap:design}.

%----------------------------------------------------------------------------------------

\section{Recruitment Software Landscape}
\label{sec:sota-recruitment}

Applicant Tracking Systems (\gls{ats}) are the dominant category of
recruitment software. At the enterprise end, suites such as Workday and
SmartRecruiters integrate requisition management, applicant tracking,
onboarding and analytics across very large organisations; mid-market tools
such as Greenhouse, Lever and Personio target structured hiring with strong
interview-workflow and collaboration features; and lighter products such as
Recruitee and Manatal aim at smaller teams, increasingly bundling
AI-assisted sourcing and ranking.

Two observations follow from this landscape and frame the present work.
First, the category is broad but its centre of gravity is \emph{workflow}
--- moving a candidate through stages --- rather than \emph{intelligence}
about the candidate pool as a persistent, analysable asset. Second, the
features that this project treats as central --- explainable, per-criterion
job matching; a deletion-aware data model built around \gls{gdpr} retention;
and an integrated language-assessment pipeline --- are typically either
absent, opaque, or offered only as separate enterprise add-ons. The platform
described here is therefore positioned not as a competitor to a full ATS but
as a right-sized talent-intelligence tool for a single multilingual HR team,
where auditability and compliance are first-class concerns.

%----------------------------------------------------------------------------------------

\section{Artificial Intelligence in Hiring}
\label{sec:sota-ai}

The recent surge of capable Large Language Models, beginning with the
demonstration that sufficiently large models perform tasks from natural
instructions with few examples \parencite{Brown2020GPT3} and continued by
efficient open foundation models \parencite{Touvron2023Llama}, has changed
what is practical in recruitment tooling. Three sub-problems are relevant
here.

\subsection{CV Parsing}
\label{sec:sota-cv}

CV parsing --- turning an unstructured document into structured fields ---
has progressed through three generations. \textbf{Rule-based} parsers
(commercial engines such as Sovren and RChilli) rely on hand-crafted
patterns and gazetteers; they are fast and predictable but brittle against
unusual layouts and multilingual content. \textbf{Machine-learning} parsers,
typically built on transformer encoders such as BERT, improve robustness by
learning to label spans, at the cost of training data and pipeline
complexity. \textbf{LLM-based} parsing, used in this project, prompts a
general model to emit a structured object directly. Its advantages are
flexibility across layouts and languages and near-zero bespoke engineering;
its risks are hallucination (invented fields), schema drift (subtly malformed
output) and per-call cost and latency. The mitigation adopted here ---
constraining the model with a strict, tolerant validation schema and a
provider fallback --- is the practical answer to those risks and is detailed
in Chapter~\ref{cap:design}.

\subsection{Candidate--Job Matching}
\label{sec:sota-matching}

Two broad approaches dominate automated matching. \textbf{Embedding-based}
ranking encodes both the candidate and the job as vectors and ranks by
similarity; it captures fuzzy semantic closeness and requires little manual
rule-writing, but it is \emph{opaque} --- a recruiter cannot easily see
\emph{why} a candidate ranked highly, which is a serious limitation where
hiring decisions must be justified and audited. \textbf{Rule-based} fit
scoring evaluates explicit criteria (skills, experience, languages,
education) and is fully explainable, at the cost of needing the criteria to
be modelled. Hybrid systems combine the two. This project deliberately
chooses the rule-based path: as argued in Chapter~\ref{cap:design}, the
ability to show a per-criterion breakdown and a clear eligibility flag was
judged more valuable for a compliance-sensitive HR team than the marginal
recall of opaque similarity.

\subsection{Conversational Assessment}
\label{sec:sota-conversational}

A growing class of products (for example HireVue, Sapia.ai and
MyInterview) conducts asynchronous, AI-mediated interviews and scores the
responses. The technical core is prompt engineering for an evaluative task:
the model must both conduct a coherent conversation and produce a structured,
rubric-aligned judgement. Such systems raise well-documented concerns around
fairness and transparency, which reinforce the design stance taken here ---
the assessment never auto-rejects a candidate; it produces an explainable,
rubric-based result for human review.

%----------------------------------------------------------------------------------------

\section{Language Proficiency Frameworks}
\label{sec:sota-cefr}

For language assessment the platform adopts the Common European Framework of
Reference for Languages (\gls{cefr}), which defines a six-level scale from A1
to C2 and describes proficiency through can-do descriptors across distinct
competences \parencite{CEFR2001}; its 2020 companion volume refines the
descriptors, notably for mediation and spoken interaction
\parencite{CEFRCompanion2020}. Compared with test-specific scales such as
TOEFL or IELTS --- which can be mapped onto CEFR bands --- CEFR is
test-agnostic, internationally recognised, and expressed as behavioural
descriptors rather than tied to a single examination. These properties make
it the natural rubric for an asynchronous, software-driven assessment: the
descriptors translate cleanly into the grammar, vocabulary and fluency
sub-scores the evaluator returns, and the resulting level is immediately
meaningful to an HR team operating across many languages.

%----------------------------------------------------------------------------------------

\section{Regulatory Context: GDPR for Recruitment}
\label{sec:sota-gdpr}

Recruitment processes personal data and therefore fall squarely under the
General Data Protection Regulation \parencite{GDPR2016}. Several of its
principles bear directly on the design. Processing requires a \emph{lawful
basis} (typically consent or legitimate interest for recruitment);
\emph{data minimisation} requires storing only what is necessary;
\emph{storage limitation} requires that data not be kept longer than needed;
and data subjects hold rights of \emph{access} and \emph{erasure}. In
practice, recruitment teams commonly translate storage limitation into a
fixed retention window --- the ``six-month rule'' that recurred throughout the
client interviews --- after which candidate data must be deleted. This single
constraint is what turns institutional memory into a liability and is the
regulatory pressure the platform is designed around: a deletion-aware data
model, an interaction-history audit trail, and a candidate self-service view
that supports the right of access, as developed in
Chapter~\ref{cap:design}.

%----------------------------------------------------------------------------------------

\section{Summary}
\label{sec:sota-summary}

The review exposes a consistent gap. Mainstream recruitment software is
strong on workflow but weak on the three properties this project treats as
essential: \emph{auditable} candidate--job matching rather than opaque
ranking; a data model built \emph{around} GDPR retention rather than bolting
compliance on afterwards; and an \emph{integrated} language-assessment
pipeline rather than a separate tool. Modern LLMs make the parsing and
assessment components feasible for a small team, provided their output is
constrained and validated, while CEFR supplies a recognised rubric for the
language dimension. The remainder of this report shows how these pieces are
combined into a single coherent product for a multilingual HR team.
